\centerline{\bf \Huge Графы Рамсея}

Кроме тривиальных чисел Рамсея, полученных в задаче 3, известно немного чисел Рамсея. Девять штук, включая те 5, которые мы рассмотрели в задачах 4 и 8. Оценки полученные там, являются точными и ниже предлагается это доказать. В других случаях рекуррентная оценка из 9 задачи оказывается грубой, и нужны дополнительно. 

Для получения нижних оценок чисел Рамсея строят примеры. Полный граф на $R(n_1, \dots, n_k) - 1$ вершинах, ребра которого раскрашены в $k$ цветов так, что для каждого цвета $i$ нет $n_i-$ подклики цвета $i$, называется \textit{графом Рамсея $R(n_1, \dots, n_k)$}. Так же предлагается понять, почему вычисление точного значения числа Рамсея алгоритмически сверхсложная задача для компьютера. 

\problem{1} Постройте граф $R(3, 3).$Докажите, что он единственный. 

\problem{2} Постройте граф $R(3, 4).$

\problem{3} Постройте граф $R(4, 4).$

\problem{4} Постройте граф $R(3, 5).$

\problem{5} Постройте граф $R(3, 3, 3).$

\problem{6} Ребра полного графа покрашены в три цвета так, что любой полный подграф на 4 вершинах содержит ребра всех цветов. Докажите, что вершин в графе не более 13. 

\problem{7} $R(3, n) \leq \dfrac{n^2 + 3}{2}$. Докажите.


$$
\alpha^3 =\left(\sqrt[3]{4+\sqrt{15}}+\sqrt[3]{4-\sqrt{15}}\right)^3 = 4 + \sqrt{15} + 4 - \sqrt{15}
$$